% ============================================================
%  Hierarchical Structure from Minimal Observations
%  A Mathematical Textbook on Constraint-Driven Entropy Reduction
%  Author: Flyxion
% ============================================================

\documentclass[11pt, twoside, openright]{book}

% ── Encoding & fonts ─────────────────────────────────────────
\usepackage[T1]{fontenc}
\usepackage[utf8]{inputenc}
\usepackage[expansion=false]{microtype}

% ── Page geometry ────────────────────────────────────────────
\usepackage[
  paperwidth=6.5in,
  paperheight=9.5in,
  top=1.1in,
  bottom=1.1in,
  inner=1.1in,
  outer=0.9in,
  marginparwidth=0pt
]{geometry}

% ── Mathematics ──────────────────────────────────────────────
\usepackage{amsmath, amssymb, amsthm, mathtools}
\usepackage{bm}
\usepackage{amsfonts}

% ── Graphics & diagrams ──────────────────────────────────────
\usepackage{tikz, pgfplots}
\usetikzlibrary{trees, arrows, positioning, shapes, shapes.geometric, calc, decorations.pathreplacing, decorations.markings}
\pgfplotsset{compat=1.18}

% ── Tables ───────────────────────────────────────────────────
\usepackage{booktabs, array, longtable}

% ── Hyperlinks ───────────────────────────────────────────────
\usepackage[
  colorlinks=true,
  linkcolor=black!70,
  citecolor=black!60,
  urlcolor=black!60
]{hyperref}

% ── Bibliography ─────────────────────────────────────────────
\usepackage[numbers, sort&compress]{natbib}

% ── Header / footer ──────────────────────────────────────────
\usepackage{fancyhdr}
\setlength{\headheight}{14pt}
\pagestyle{fancy}
\fancyhf{}
\fancyhead[LE]{\small\itshape\leftmark}
\fancyhead[RO]{\small\itshape\rightmark}
\fancyfoot[C]{\small\thepage}
\renewcommand{\headrulewidth}{0.4pt}

% ── Chapter titles ───────────────────────────────────────────
\usepackage{titlesec}
\titleformat{\chapter}[display]
  {\normalfont\Large\bfseries}
  {\chaptertitlename\ \thechapter}{16pt}{\LARGE}

% ── Theorem environments ─────────────────────────────────────
\theoremstyle{plain}
\newtheorem{theorem}{Theorem}[chapter]
\newtheorem{lemma}[theorem]{Lemma}
\newtheorem{proposition}[theorem]{Proposition}
\newtheorem{corollary}[theorem]{Corollary}

\theoremstyle{definition}
\newtheorem{definition}[theorem]{Definition}
\newtheorem{example}[theorem]{Example}
\newtheorem{algorithm}[theorem]{Algorithm}

\theoremstyle{remark}
\newtheorem{remark}[theorem]{Remark}
\newtheorem{notation}[theorem]{Notation}

% ── Diagram styles ───────────────────────────────────────────
\usepackage{xcolor}
\definecolor{nodeblue}{RGB}{70,130,180}
\definecolor{nodered}{RGB}{180,70,70}
\definecolor{nodegreen}{RGB}{60,150,90}
\definecolor{nodegray}{RGB}{140,140,140}
\definecolor{supercolor}{RGB}{210,140,50}
\definecolor{constraintblue}{RGB}{100,160,220}

\tikzset{
  treenode/.style={circle, draw=nodeblue!80, fill=nodeblue!15,
                   minimum size=7mm, inner sep=1pt, font=\small},
  leafnode/.style={circle, draw=nodegreen!80, fill=nodegreen!20,
                   minimum size=7mm, inner sep=1pt, font=\small},
  internalnode/.style={circle, draw=nodeblue!80, fill=nodeblue!20,
                   minimum size=6mm, inner sep=1pt, font=\tiny},
  supervertex/.style={rectangle, rounded corners=4pt, draw=supercolor!80,
                   fill=supercolor!20, minimum size=8mm, inner sep=3pt, font=\small},
  ghostnode/.style={circle, draw=nodegray!60, fill=nodegray!10,
                   minimum size=6mm, inner sep=1pt, font=\tiny, dashed},
  treeedge/.style={draw=nodeblue!60, thick},
  highlightedge/.style={draw=nodered!80, thick},
  constraint/.style={draw=constraintblue!80, dashed, thick},
  selectionfilter/.style={trapezium, trapezium left angle=70,
                   trapezium right angle=110, draw=nodered!70,
                   fill=nodered!10, minimum height=8mm, font=\small}
}

% ── Custom commands ──────────────────────────────────────────
\newcommand{\LCA}{\operatorname{LCA}}
\newcommand{\cost}{\operatorname{cost}}
\newcommand{\reward}{\operatorname{reward}}
\newcommand{\leaves}{\operatorname{leaves}}
\newcommand{\FDR}{\mathrm{FDR}}
\newcommand{\poly}{\operatorname{poly}}
\newcommand{\polylog}{\operatorname{polylog}}
\newcommand{\Obb}{\mathcal{O}}
\newcommand{\Hyp}{\mathcal{H}}
\newcommand{\Tree}{\mathcal{T}}
\newcommand{\Cset}{\mathcal{C}}
\newcommand{\Rset}{\mathcal{R}}
\newcommand{\KL}{\mathrm{KL}}
\newcommand{\dd}{\,\mathrm{d}}
\newcommand{\eps}{\varepsilon}
\newcommand{\wh}[1]{\widehat{#1}}
\newcommand{\norm}[1]{\left\|#1\right\|}
\DeclareMathOperator*{\argmin}{arg\,min}
\DeclareMathOperator*{\argmax}{arg\,max}

% ── Front matter / back matter helpers ───────────────────────
\usepackage{tocloft}
\setlength{\cftbeforechapskip}{6pt}

% ── Epigraph ─────────────────────────────────────────────────
\usepackage{epigraph}
\setlength{\epigraphwidth}{0.62\textwidth}

% ── Index ────────────────────────────────────────────────────
\usepackage{makeidx}
\makeindex

% ============================================================
\begin{document}

\frontmatter
\include{chapters/titlepage}
\include{chapters/preface}
\tableofcontents
\listoffigures

\mainmatter

%% ── PART I ──────────────────────────────────────────────────
\part{Prerequisites and Foundations}
\include{chapters/ch01_orientation}
\include{chapters/ch02_discrete}
\include{chapters/ch03_probability}
\include{chapters/ch04_entropy}
\include{chapters/ch05_metrics}
\include{chapters/ch06_constraints}

%% ── PART II ─────────────────────────────────────────────────
\part{Hierarchical Structure and Minimal Relational Information}
\include{chapters/ch07_dendrograms}
\include{chapters/ch08_ultrametric}
\include{chapters/ch09_triplets}
\include{chapters/ch10_reconstruction}

%% ── PART III ────────────────────────────────────────────────
\part{Classical Objectives and Hardness}
\include{chapters/ch11_objectives}
\include{chapters/ch12_hardness}

%% ── PART IV ─────────────────────────────────────────────────
\part{Learning-Augmented Clustering and the Splitting Oracle}
\include{chapters/ch13_oracle}
\include{chapters/ch14_partial_hc}
\include{chapters/ch15_topdown}

%% ── PART V ──────────────────────────────────────────────────
\part{Approximation Guarantees and Algorithmic Consequences}
\include{chapters/ch16_approx}
\include{chapters/ch17_complexity}

%% ── PART VI ─────────────────────────────────────────────────
\part{Streaming, Parallelism, and Compressed Inference}
\include{chapters/ch18_streaming}
\include{chapters/ch19_parallel}

%% ── PART VII ────────────────────────────────────────────────
\part{Geometric and Information-Theoretic Reinterpretations}
\include{chapters/ch20_curvature}
\include{chapters/ch21_entropy_descent}
\include{chapters/ch22_constraint_prop}

%% ── PART VIII ───────────────────────────────────────────────
\part{Inference Systems, Replication, and Knowledge Accumulation}
\include{chapters/ch23_bayesian}
\include{chapters/ch24_replication}
\include{chapters/ch25_operators}
\include{chapters/ch26_exploratory}

%% ── PART IX (Appendices) ────────────────────────────────────
\part{Full Derivations and Advanced Appendices}
\appendix
\include{chapters/appA_ultrametric_equiv}
\include{chapters/appB_triplet_topology}
\include{chapters/appC_noise_analysis}
\include{chapters/appD_collapse_loss}
\include{chapters/appE_variational}

\backmatter
\include{chapters/synthesis}
\bibliographystyle{plainnat}
\bibliography{references}
\printindex

\end{document}
